% Motor: XeLaTeX. Sesión 21 - material docente UTP.
% Estructura docente de la Sesión 20, condensada a 20 diapositivas.
\documentclass[aspectratio=169,10pt,c]{beamer}
% Reutiliza sin cambios la identidad visual UTP validada en la Sesion 08.
% Reutiliza sin cambios la identidad visual UTP validada en la Sesion 05.
% La segunda ruta permite compartir los assets comunes sin duplicar binarios.
% !TEX root = slides.tex
% Preambulo comun: estilo UTP, colores, fuentes, plantillas de beamer y
% comandos de apoyo didactico. Se carga desde slides.tex tras \documentclass.

% La opción "c" de beamer reparte el espacio libre 1:1.5 (arriba:abajo), por lo
% que el contenido queda por encima del centro. Igualamos ambos rellenos.
\makeatletter
\define@key{beamerframe}{c}[true]{%
  \beamer@frametopskip=0pt plus 1fill\relax%
  \beamer@framebottomskip=0pt plus 1fill\relax%
  \beamer@frametopskipautobreak=0pt plus .5\paperheight\relax%
  \beamer@framebottomskipautobreak=0pt plus .5\paperheight\relax%
}
\makeatother

\usepackage{graphicx}
\usepackage{tikz}
\usepackage{xcolor}
\usepackage{iftex}

\ifPDFTeX
  \usepackage[utf8]{inputenc}
  \usepackage[T1]{fontenc}
  \usepackage{lmodern}          % tipografia escalable para el respaldo con pdflatex
  \renewcommand{\familydefault}{\sfdefault}
\else
  \usepackage{fontspec}
  \IfFontExistsTF{Arial}
    {\setsansfont{Arial}[BoldFont={Arial Bold}]}
    {\setsansfont{Helvetica}[BoldFont={Helvetica Bold}]}
\fi

\usepackage[spanish,es-nodecimaldot]{babel}
\usepackage{microtype}
\usepackage{booktabs}
\usepackage{array}
\usetikzlibrary{arrows.meta,positioning,fit,calc,shapes.geometric}

\graphicspath{{assets/}}

% ---------- Identidad visual UTP ----------
\definecolor{utpRojo}{HTML}{C42B35}
\definecolor{utpNegro}{HTML}{111111}
\definecolor{utpGris}{HTML}{5F5F5F}
\definecolor{utpGrisMedio}{HTML}{9A9A9A}
\definecolor{utpGrisClaro}{HTML}{F1F1F1}

\usetheme{default}
\setbeamercolor{normal text}{fg=utpNegro,bg=white}
\setbeamercolor{structure}{fg=utpRojo}
\setbeamercolor{frametitle}{fg=utpNegro,bg=white}
\setbeamercolor{title}{fg=utpNegro}
\setbeamercolor{subtitle}{fg=utpNegro}
\setbeamercolor{author}{fg=utpNegro}
\setbeamercolor{date}{fg=utpNegro}
\setbeamercolor{block title}{fg=white,bg=utpNegro}
\setbeamercolor{block body}{fg=utpNegro,bg=utpGrisClaro}
\setbeamercolor{alertblock title}{fg=white,bg=utpRojo}
\setbeamercolor{alertblock body}{fg=utpNegro,bg=utpRojo!7}
\setbeamercolor{exampleblock title}{fg=white,bg=utpGris}
\setbeamercolor{exampleblock body}{fg=utpNegro,bg=utpGris!8}
\setbeamercolor{terminal}{fg=white,bg=utpNegro}

\setbeamerfont{title}{series=\bfseries,size=\fontsize{19.5}{22.5}\selectfont}
\setbeamerfont{subtitle}{size=\fontsize{11.5}{13.5}\selectfont}
\setbeamerfont{date}{series=\bfseries,size=\fontsize{11.5}{13.5}\selectfont}
\setbeamerfont{author}{size=\fontsize{10.5}{12.5}\selectfont}
\setbeamerfont{frametitle}{series=\bfseries,size=\large}
\setbeamerfont{block title}{series=\bfseries,size=\small}
\setbeamerfont{block body}{size=\normalsize}

\setbeamertemplate{navigation symbols}{}
\setbeamertemplate{headline}{}
\setbeamertemplate{blocks}[default]
\setbeamertemplate{itemize item}{\color{utpRojo}\small\raise.15ex\hbox{$\blacktriangleright$}}
\setbeamertemplate{itemize subitem}{\color{utpGris}\scriptsize$\bullet$}
\setbeamertemplate{enumerate item}{\color{utpRojo}\insertenumlabel.}
\setbeamertemplate{background canvas}{\color{white}\rule{\paperwidth}{\paperheight}}

% Título de diapositiva: negrita negra a la izquierda, logo UTP a la derecha
\setbeamertemplate{frametitle}{%
  \begin{tikzpicture}[remember picture,overlay]
    \node[anchor=north east,inner sep=0pt]
      at ([xshift=-7mm,yshift=-5.5mm]current page.north east)
      {\includegraphics[width=17mm]{utp-logo.png}};
  \end{tikzpicture}%
  \vspace{1.2mm}%
  \begin{minipage}[t]{.80\textwidth}\raggedright\insertframetitle\end{minipage}\par
  \vspace{1.1mm}{\color{utpRojo}\rule{\textwidth}{1.1pt}}
}

% Solo numeración discreta; la plantilla UTP no lleva pie de página
\setbeamertemplate{footline}{%
  \hfill{\color{utpGrisMedio}\scriptsize\insertframenumber}\hspace{6mm}\vspace{2.5mm}%
}

% Portada UTP
\setbeamertemplate{title page}{%
  \begin{tikzpicture}[remember picture,overlay]
    \node[anchor=center,inner sep=0pt] at (current page.center) {%
      \includegraphics[width=\paperwidth,height=\paperheight]{utp-cover-bg.jpeg}%
    };
    \node[anchor=north,align=center,text=utpNegro,inner sep=0pt,
      text width=145mm,font=\usebeamerfont{title}]
      at ([yshift=-25.2mm]current page.north) {\inserttitle};
    \node[anchor=north,align=center,text=utpNegro,inner sep=0pt,
      text width=142mm,font=\usebeamerfont{subtitle}]
      at ([yshift=-38.2mm]current page.north) {\insertsubtitle};
    \node[anchor=north,align=center,text=utpNegro,inner sep=0pt,
      font=\usebeamerfont{date}]
      at ([yshift=-50.4mm]current page.north) {\insertdate};
    \node[anchor=north,align=center,text=utpNegro,inner sep=0pt,
      font=\usebeamerfont{author}]
      at ([yshift=-58.4mm]current page.north) {\insertauthor};
  \end{tikzpicture}%
}

% ---------- Comandos de apoyo didáctico ----------
\newcommand{\concepto}[2]{%
  \begin{block}{#1}\large #2\end{block}}
\newcommand{\pregunta}[1]{%
  \begin{alertblock}{Decide y explica}#1\end{alertblock}}
\newcommand{\analogia}[1]{%
  \begin{exampleblock}{En palabras simples}\small #1\end{exampleblock}}
\newcommand{\enlarealidad}[1]{%
  \begin{block}{Cómo se ve en la realidad}\small #1\end{block}}
\newcommand{\ojo}[1]{%
  \begin{alertblock}{Ojo con esto}\small #1\end{alertblock}}
\newcommand{\apoyo}{{\normalfont\footnotesize\color{utpGris}~· material de apoyo}}
\newcommand{\marca}[1]{\tikz[baseline=(m.base)]{%
  \node[rounded corners=1.5pt,fill=utpGrisClaro,draw=utpGrisMedio,inner xsep=3pt,inner ysep=1.5pt,
        font=\footnotesize] (m) {#1};}}
\newenvironment{terminal}{%
  \begin{beamercolorbox}[wd=\linewidth,sep=5pt,rounded=true]{terminal}%
  \ttfamily\small\ignorespaces}{%
  \end{beamercolorbox}}

\tikzset{
  caja/.style={draw=utpNegro, very thick, rounded corners=1.5pt, fill=white,
    align=center, minimum height=0.72cm, inner sep=5pt, font=\small},
  cajaazul/.style={caja, draw=utpGris, fill=white},
  cajaverde/.style={caja, draw=utpNegro, fill=utpGrisClaro},
  cajanaranja/.style={caja, draw=utpRojo!75, fill=utpRojo!7},
  flecha/.style={-{Stealth[length=2.3mm]}, very thick, draw=utpNegro},
  flechavuelta/.style={{Stealth[length=2.3mm]}-{Stealth[length=2.3mm]}, very thick, draw=utpNegro}
}

\graphicspath{{../../Sesion05/slides/assets/}}



\graphicspath{{assets/}{../../Sesion05/slides/assets/}}

\title{Integración de Servicios Cloud GCP}
\subtitle{De una idea a una webapp: frontend, backend, datos y archivos}
\date{Semana 11 - Sesión 21}
\author{Rel Guzman}

\newcommand{\separador}[2]{%
\begin{frame}[plain]
\vfill
\begin{center}
{\fontsize{29}{34}\selectfont\bfseries #1}\par
\vspace{2mm}
{\small #2}\par
\vspace{4mm}
{\color{utpRojo}\rule{8.5cm}{1.2pt}}
\end{center}
\vfill
\end{frame}}

\begin{document}

{\setbeamertemplate{footline}{}\maketitle}

% 2
\begin{frame}{Revisión: los bloques que ya construiste}
\renewcommand{\arraystretch}{1.28}
\centering\small
\begin{tabular}{@{}p{3.0cm}p{5.0cm}p{5.0cm}@{}}
\toprule
\textbf{Laboratorio} & \textbf{Lo que hiciste} & \textbf{El bloque de una webapp}\\
\midrule
Lab 01 & Creaste buckets, subiste objetos y usaste Cloud Shell & \textbf{Archivos} y una forma de administrar recursos\\
Lab 02 & Implementaste una pila LAMP desde Marketplace & \textbf{Servidor web, PHP y base de datos} en una VM\\
Lab 03 y 04 & Probaste VPC, una VM privada, Acceso privado y Cloud NAT & \textbf{Red} y salida segura de los servidores\\
Lab 05 y 06 & Configuraste VMs, discos, firewall, copias y scripts & \textbf{Cómputo}, datos persistentes y operación\\
\bottomrule
\end{tabular}
\par\vspace{3mm}
\begin{alertblock}{La pregunta de hoy}
\small Cuando una aplicación necesita varios bloques, ¿cuál va en cada lugar y por qué?
\end{alertblock}
\end{frame}

% 3
\begin{frame}{Revisión: una aplicación que sí terminaste}
\begin{block}{El servidor de Minecraft del Lab 06}
\small No era solo una VM. Tenía una IP externa para los jugadores, una regla de firewall para el puerto \texttt{25565}, un disco para el mundo, copias en un bucket y scripts para iniciar y detener el servicio.
\end{block}
\begin{center}
\begin{tikzpicture}[font=\small,
 caja/.style={draw=utpRojo,rounded corners=4pt,minimum width=2.55cm,minimum height=1.0cm,align=center}]
  \node[caja,fill=utpRojo!25] (u) at (0,0) {\textbf{Jugadores}\\\footnotesize Internet};
  \node[caja,fill=white,very thick] (v) at (4,0) {\textbf{VM}\\\footnotesize Minecraft};
  \node[caja,fill=utpRojo!10] (d) at (8,0.8) {\textbf{Disco}\\\footnotesize mundo};
  \node[caja,fill=utpRojo!10] (b) at (8,-0.8) {\textbf{Bucket}\\\footnotesize copias};
  \draw[-latex,thick] (u) -- node[above,font=\footnotesize] {firewall} (v);
  \draw[-latex,thick] (v) -- (d);
  \draw[-latex,thick] (v) -- (b);
\end{tikzpicture}
\end{center}
\begin{alertblock}{Pregunta rápida}
\small Si guardas el mundo solo dentro de la VM, ¿qué riesgo corres cuando la borras o su disco falla?
\end{alertblock}
\end{frame}

% 4
\begin{frame}[c]
\begin{block}{Logro de la sesión}
Al finalizar, podrás reconocer las partes básicas de una webapp, explicar qué hace cada una y proponer una arquitectura inicial en Google Cloud que conecte \textbf{frontend}, \textbf{backend}, \textbf{base de datos}, \textbf{archivos} y \textbf{seguridad}.
\end{block}
\end{frame}

% 5
\begin{frame}{Saberes previos}
\begin{itemize}\setlength\itemsep{11pt}
  \item Cuando abres una tienda en línea, ¿qué ves primero: el sitio, la base de datos o el servidor?
  \item Si una persona sube una foto de perfil, ¿debería quedar en la base de datos o como archivo?
  \item En una aplicación con usuarios y pedidos, ¿qué información no debería ser pública?
\end{itemize}
\end{frame}

% 6
\begin{frame}[c]{Utilidad de esta sesión}
\begin{center}
\begin{minipage}{.84\textwidth}
{\large En un proyecto real no eliges un servicio por su nombre. Primero identificas \textbf{qué necesita la aplicación}: mostrar una página, ejecutar una regla, guardar datos, conservar archivos o proteger una clave. Luego eliges el servicio que cumple esa función.}
\end{minipage}
\end{center}
\end{frame}

% 7
\begin{frame}{Plan de hoy: dos partes y una práctica}
\renewcommand{\arraystretch}{1.35}
\centering\small
\begin{tabular}{@{}p{1.55cm}>{\raggedright\arraybackslash}p{5.8cm}>{\raggedright\arraybackslash}p{5.45cm}@{}}
\toprule
\textbf{Parte} & \textbf{Qué vemos} & \textbf{Lo conectamos con}\\
\midrule
1. Teoría & Las cinco piezas de una webapp & La tienda, la galería y el servidor de Minecraft\\
2. Teoría & Cuándo usar un servicio gestionado o una VM & Lo que ya configuraste en los labs\\
\textbf{Práctica} & Diseñar la arquitectura inicial de un inventario & El caso del curso y tus decisiones\\
\bottomrule
\end{tabular}
\par\vspace{3mm}
\begin{alertblock}{Regla de hoy}
\small Una aplicación no tiene que usar todos los servicios. Solo usa los bloques que realmente necesita.
\end{alertblock}
\end{frame}

% 8
\separador{Parte 1: las piezas de una webapp}{Del navegador a los datos y archivos}

% 9
\begin{frame}{Una webapp en cinco piezas}
\begin{center}
\begin{tikzpicture}[font=\small,
 caja/.style={draw=utpRojo,rounded corners=4pt,minimum width=2.45cm,minimum height=1.0cm,align=center}]
  \node[caja,fill=utpRojo!25] (user) at (0,0) {\textbf{Persona}\\\footnotesize navegador};
  \node[caja,fill=utpRojo!10] (front) at (3.2,0) {\textbf{Frontend}\\\footnotesize lo que ve};
  \node[caja,fill=white,very thick] (back) at (6.4,0) {\textbf{Backend}\\\footnotesize las reglas};
  \node[caja,fill=utpRojo!10] (db) at (9.6,0.9) {\textbf{Base de datos}\\\footnotesize usuarios y pedidos};
  \node[caja,fill=utpRojo!10] (files) at (9.6,-0.9) {\textbf{Almacenamiento}\\\footnotesize fotos y documentos};
  \draw[-latex,thick] (user) -- (front);
  \draw[-latex,thick] (front) -- (back);
  \draw[-latex,thick] (back) -- (db);
  \draw[-latex,thick] (back) -- (files);
\end{tikzpicture}
\end{center}
\renewcommand{\arraystretch}{1.24}
\centering\small
\begin{tabular}{@{}p{2.8cm}p{8.9cm}@{}}
\toprule
\textbf{Pieza} & \textbf{Su trabajo}\\
\midrule
Frontend & Muestra páginas, botones e imágenes en el navegador.\\
Backend & Recibe pedidos, valida reglas y conversa con los demás servicios.\\
Datos y archivos & Guardan información que debe sobrevivir cuando la aplicación se reinicia.\\
\bottomrule
\end{tabular}
\end{frame}

% 10
\begin{frame}{Frontend: lo que llega al navegador}
\begin{columns}[T,onlytextwidth]
\column{.48\textwidth}
\begin{block}{Qué contiene}
\small HTML, CSS, JavaScript, imágenes, íconos y fuentes. Es la parte visible de una landing, un catálogo o un panel.
\end{block}
\begin{exampleblock}{Caso simple}
\small Si las páginas ya están listas antes de la visita, puedes publicarlas como sitio estático. El navegador descarga los archivos y los muestra.
\end{exampleblock}
\column{.48\textwidth}
\begin{block}{En Google Cloud}
\small Un bucket puede guardar objetos. Para un sitio web público debes decidir explícitamente qué contenido será público y qué contenido se mantendrá privado.
\end{block}
\begin{alertblock}{Recuerda el Lab 01}
\small Ya creaste un bucket y copiaste archivos desde Cloud Shell. El mismo bloque sirve para conservar archivos de una aplicación.
\end{alertblock}
\end{columns}
\end{frame}

% 11
\begin{frame}{Backend: donde viven las reglas de la aplicación}
\begin{block}{Qué hace}
\small Una API o backend recibe una petición, comprueba permisos, aplica una regla y devuelve una respuesta. Por ejemplo: ``registrar pedido'', ``calcular total'' o ``mostrar mis facturas''.
\end{block}
\begin{columns}[T,onlytextwidth]
\column{.48\textwidth}
\begin{exampleblock}{Cuando una VM tiene sentido}
\small Necesitas controlar el sistema operativo, instalar software propio o mantener un proceso encendido. El servidor de Minecraft y la pila LAMP fueron ejemplos de este caso.
\end{exampleblock}
\column{.48\textwidth}
\begin{exampleblock}{Cuando conviene que lo gestione el proveedor}
\small Para una API pequeña o un contenedor, un servicio gestionado evita administrar parches, arranque y escalado de servidores.
\end{exampleblock}
\end{columns}
\par\vspace{2mm}
\begin{alertblock}{La decisión no cambia la función}
\small En ambos casos el backend sigue siendo el lugar donde viven las reglas. Lo que cambia es quién administra la infraestructura.
\end{alertblock}
\end{frame}

% 12
\begin{frame}{Base de datos: información que la aplicación consulta}
\renewcommand{\arraystretch}{1.32}
\centering\small
\begin{tabular}{@{}p{3.3cm}p{4.7cm}p{4.3cm}@{}}
\toprule
\textbf{Si guardas\ldots} & \textbf{Una opción básica} & \textbf{Ejemplo}\\
\midrule
Tablas relacionadas & Base de datos relacional o SQL & Usuarios, productos, pedidos y facturas\\
Documentos flexibles & Base de datos de documentos o NoSQL & Perfiles con campos que cambian según el usuario\\
Datos temporales & Caché de clave y valor & Sesiones, contadores o resultados rápidos\\
\bottomrule
\end{tabular}
\par\vspace{3mm}
\begin{columns}[T,onlytextwidth]
\column{.48\textwidth}
\begin{block}{No confundas dato con archivo}
\small Un pedido puede guardar el nombre de una factura y su enlace. El PDF completo es un archivo, no una fila de texto.
\end{block}
\column{.48\textwidth}
\begin{alertblock}{Conexión con el Lab 02}
\small MariaDB de LAMP era una base relacional dentro de la VM. Una base gestionada mantiene la misma función, pero el proveedor administra el servidor.
\end{alertblock}
\end{columns}
\end{frame}

% 13
\begin{frame}{Archivos y secretos: no todo se publica}
\renewcommand{\arraystretch}{1.35}
\centering\small
\begin{tabular}{@{}p{3.3cm}p{4.1cm}p{4.9cm}@{}}
\toprule
\textbf{Recurso} & \textbf{Dónde va} & \textbf{Regla básica}\\
\midrule
Foto de un producto & Bucket público o entrega pública & Cualquiera puede verla si forma parte del catálogo\\
Factura o CV & Bucket privado y enlace temporal & El backend revisa quién puede descargarlo\\
Copia de seguridad y registros & Bucket privado separado & No se exponen desde la aplicación\\
Clave de API o archivo \texttt{.env} & Gestor de secretos & Nunca en el repositorio ni en un bucket\\
\bottomrule
\end{tabular}
\par\vspace{3mm}
\begin{alertblock}{La idea del Lab 06}
\small El mundo de Minecraft vivía en un disco y sus copias iban a un bucket. Separar la aplicación de su copia reduce el riesgo de perder ambos a la vez.
\end{alertblock}
\end{frame}

% 14
\separador{Parte 2: elegir dónde corre cada pieza}{Servicio gestionado, VM y seguridad}

% 15
\begin{frame}{Servicio gestionado o VM}
\renewcommand{\arraystretch}{1.35}
\centering\small
\begin{tabular}{@{}p{3.2cm}p{4.45cm}p{4.75cm}@{}}
\toprule
\textbf{Pregunta} & \textbf{Servicio gestionado} & \textbf{Máquina virtual}\\
\midrule
¿Quién aplica parches? & El proveedor se encarga del servicio & Tú actualizas el sistema operativo y el software\\
¿Cuándo encaja? & Necesitas una función común: API, contenedor, base de datos o cola & Necesitas control del sistema, software especial o un proceso permanente\\
¿Qué debes vigilar? & Configuración, permisos y costos & Lo anterior, más disco, firewall, copias y mantenimiento\\
\bottomrule
\end{tabular}
\par\vspace{3mm}
\begin{exampleblock}{No hay una respuesta única}
\small Una webapp puede usar frontend y base de datos gestionados, pero conservar una VM para un servicio que debe estar encendido todo el tiempo.
\end{exampleblock}
\end{frame}

% 16
\begin{frame}{Una arquitectura inicial en Google Cloud}
\begin{center}
\begin{tikzpicture}[font=\small,
 caja/.style={draw=utpRojo,rounded corners=4pt,minimum width=2.55cm,minimum height=1.0cm,align=center}]
  \node[caja,fill=utpRojo!25] (u) at (0,0) {\textbf{Cliente}\\\footnotesize navegador};
  \node[caja,fill=utpRojo!10] (f) at (3.0,0) {\textbf{Frontend}\\\footnotesize archivos públicos};
  \node[caja,fill=white,very thick] (a) at (6.0,0) {\textbf{Backend}\\\footnotesize VM o servicio};
  \node[caja,fill=utpRojo!10] (d) at (9.2,0.85) {\textbf{BD}\\\footnotesize datos privados};
  \node[caja,fill=utpRojo!10] (s) at (9.2,-0.85) {\textbf{Bucket}\\\footnotesize archivos privados};
  \draw[-latex,thick] (u) -- (f);
  \draw[-latex,thick] (f) -- (a);
  \draw[-latex,thick] (a) -- (d);
  \draw[-latex,thick] (a) -- (s);
\end{tikzpicture}
\end{center}
\begin{columns}[T,onlytextwidth]
\column{.48\textwidth}
\begin{block}{Lo que ya sabes proteger}
\small La VPC organiza la red; el firewall decide qué tráfico entra; Cloud NAT permite que una VM privada salga a Internet para actualizarse.
\end{block}
\column{.48\textwidth}
\begin{block}{Lo que no debe cambiar}
\small El backend debe conservar acceso solo a los datos y archivos que necesita. Las claves van en un gestor de secretos.
\end{block}
\end{columns}
\end{frame}

% 17
\begin{frame}{Caso: inventario de una tienda}
\begin{block}{La necesidad}
\small La tienda muestra productos, permite iniciar sesión, genera reportes, guarda facturas de proveedores y necesita respaldos. El equipo también quiere conservar las claves fuera del código.
\end{block}
\renewcommand{\arraystretch}{1.26}
\centering\small
\begin{tabular}{@{}p{4.0cm}p{5.2cm}p{3.3cm}@{}}
\toprule
\textbf{Elemento} & \textbf{Decisión inicial} & \textbf{Razón}\\
\midrule
Catálogo y fotos de producto & Frontend y archivos de lectura pública & Los clientes deben verlos\\
Usuarios, productos y pedidos & Base de datos relacional & Hay relaciones entre las tablas\\
Facturas y reportes & Bucket privado con entrega controlada & Solo personal autorizado los consulta\\
Claves de producción & Gestor de secretos & No viajan con el código\\
Respaldos y registros & Bucket privado separado & No forman parte del sitio público\\
\bottomrule
\end{tabular}
\end{frame}

% 18
\begin{frame}{Práctica: dibuja una arquitectura inicial}
\begin{block}{En grupos de dos}
\small Elijan una galería de fotos, una plataforma de entregas o una biblioteca de grabaciones. Dibujen solo los bloques necesarios y las flechas importantes.
\end{block}
\renewcommand{\arraystretch}{1.15}
\centering\small
\begin{tabular}{@{}p{2.8cm}p{6.2cm}p{3.2cm}@{}}
\toprule
\textbf{Decidan} & \textbf{Pregunta que deben responder} & \textbf{Evidencia de los labs}\\
\midrule
Frontend & ¿Qué ve el visitante? & Bucket y objetos del Lab 01\\
Backend & ¿Qué regla ejecuta la aplicación? & LAMP o VM de los Labs 02, 05 y 06\\
Datos y archivos & ¿Qué queda privado? ¿Qué necesita copia? & Disco, bucket y backup del Lab 06\\
Red y acceso & ¿Qué puerto se abre? ¿Qué queda privado? & VPC, firewall y NAT de los Labs 03 y 04\\
\bottomrule
\end{tabular}
\par\vspace{0.5mm}
\begin{alertblock}{Entregable}
\small Un diagrama con al menos cuatro bloques y una frase que justifique cada decisión importante.
\end{alertblock}
\end{frame}

% 19
\begin{frame}{Cierre: preguntas para comprobar}
\begin{block}{Responde con tus palabras}
\small
1. ¿Qué hace el frontend y qué hace el backend?\\\medskip
2. ¿Por qué una factura PDF no debe guardarse igual que un pedido?\\\medskip
3. ¿Qué diferencia hay entre un bucket público y uno privado?\\\medskip
4. ¿Cuándo una VM tiene más sentido que un servicio gestionado?\\\medskip
5. En la tienda, ¿qué elementos dejarías sin IP externa y por qué?\\\medskip
6. ¿Qué laboratorio usarías como ejemplo para explicar copias de seguridad?
\end{block}
\begin{alertblock}{Idea final}
\small Una arquitectura comienza al separar responsabilidades: interfaz, reglas, datos, archivos y seguridad.
\end{alertblock}
\end{frame}

% 20
\begin{frame}{Referencias}
\scriptsize
\begin{itemize}\setlength\itemsep{4pt}
  \item Material base de esta sesión: \texttt{/Users/rgap/Desktop/Notas\_Clases\_GCP}: diagramas de frontend, API/backend, bases de datos, almacenamiento y otras cargas.
  \item Casos de práctica: \texttt{Notas\_Clases\_GCP/niveles/niveles.md}.
  \item Evidencias de laboratorios: \texttt{/Users/rgap/Desktop/Modelos\_labs/Lab01} a \texttt{Lab06}.
  \item Material previo del curso: Sesiones 19 y 20, redes VPC, Acceso privado a Google y Cloud NAT.
  \item Google Cloud: Cloud Storage, Compute Engine, VPC, Cloud NAT, Cloud SQL, Firestore y Secret Manager.
  \item Las respuestas de las preguntas y de la práctica están en \texttt{../respuestas.md}.
\end{itemize}
\end{frame}

\end{document}
